

Or ces
deux carrés sont les carrés «~de changement de base~» déduits des deux carrés (7.10.0), de sorte qu'on peut appliquer (IV 3). Il faut seulement vérifier que ces carrés correspondent à des couples \((X,Y_t)\) et \((X,Y_s)\) de schémas sur \(Y\) qui sont tor-indépendants sur \(Y\). Cela est trivial pour le premier couple, \(Y_t\to Y\) étant plat, et aussi immédiat pour le deuxième, l'immersion \(Y_s\to Y\) étant régulière de codimension 1 et restant ainsi après le changement de base \(X\to Y\).

Ne supposons plus maintenant \(g\) propre, mais supposons-le en revanche de tor-dimension finie. Je dis qu'il en est alors de même de \(g_t\) et \(g_s\). C'est trivial pour le premier, et pour le second résulte du fait (qu'on vient de signaler) que \(X\) et \(Y_s\) sont tor-indépendants sur \(Y\). On peut donc considérer le carré d'homomorphismes
\begin{equation}\tag{7.10.2}
\begin{tikzcd}[column sep=large,row sep=large]
K_\bullet(Y_t) \arrow[r,"g_t^*"] \arrow[d,"\sigma_Y"'] &
K_\bullet(X_t) \arrow[d,"\sigma_X"]\\
K_\bullet(Y_s) \arrow[r,"g_s^*"] &
K_\bullet(X_s),
\end{tikzcd}
\end{equation}
où les flèches horizontales sont définies en (IV 2.12). On notera que si \(X\) et \(Y\) satisfont aux conditions de 7.7 i.e. sont à fibres régulières, alors (\(Y\) étant régulier) \(g\) est automatiquement de tor-dimension finie, et le diagramme précédent peut s'interpréter comme la commutativité d'un diagramme d'anneaux \(K^\bullet\).

Pour prouver la commutativité de (7.10.2), il suffit encore de prouver celle des carrés antérieurs du prisme analogue à \((*)\), mais où les homomorphismes image directe \(g_{t*}\), \(g_*\), \(g_{s*}\) sont remplacés par les homomorphismes image inverse \(g_t^*\), \(g^*\), \(g_s^*\). Or la commutativité de ces carrés résulte aussi%
